\section{Conclusion}

This thesis examined the following questions: Does information about the relationship between Bitcoin and conventional assets help predict how that relationship will change, and can these predictions aid in risk management? The answer to both questions is yes. However, it is necessary to clarify what ``useful'' means.

The relationship between assets can be predicted based on a sample. This conclusion holds true for six asset pairs, four different time intervals, and several model types tested using a step-by-step procedure. Longer intervals make the trend smoother and more stable over time, leading to more accurate forecasts. With a 90-day interval, even a simple AR(1) model yields an $R^2$ ranging from 0.982 to 0.992 across all six pairs.

This finding also points to an important limitation: the forecast is largely based on constancy. Since the 30-day moving correlation matches the previous day's data in 29 out of 30 observations, it is difficult to make a forecast that would be any better than simply saying that tomorrow's correlation will be the same as today's. The Ridge, AR(1), and HAR models fall within the range of statistical noise, and the simple last-value forecast is just behind them. The Random Forest, GBM, and XGBoost models perform worse. Tree-based models do not outperform the baseline autoregressive model, while Ridge performs at a similar level. All machine learning models perform better than the DCC-GARCH test.

This last comparison warrants careful consideration. The target metric is the moving Pearson correlation, whereas DCC-GARCH estimates the current conditional correlation. Comparing the DCC-GARCH output to the target metric introduces a lag that does not align with the benchmark. Thus, the main takeaway is limited: flexible models fit this specific objective more accurately than a parametric model, but they do not perform better than a simple stationary model. This does not mean that DCC-GARCH generally performs worse than a market-dependent model.

The removal of feature groups in Section 4.6 illustrates the same idea from another perspective. A single lag in the dependent variable already accounts for most of the variation that can be predicted. Adding more historical data to the dependent variable yields only a slight improvement. Adding groups of interrelated objects does not improve the linear model and actually worsens the tree-based models. Tree-based models perform worse as more groups of objects are added.

The investor signal level is more complex. The balanced accuracy ranges from 0.50 to 0.53, while the AUC ranges from 0.48 to 0.56. These values are modest. The probability of exceeding them is high for most, but not all, configurations. The probability of a decline in the Bitcoin–NASDAQ pair is lower over longer time frames. The strongest signal is observed for the Bitcoin–S\&P 500 pair, where the characteristics of the cryptocurrency are more closely linked to factors affecting the broader stock market. For gold, silver, and the dollar index, the signal is weaker. This weak signal makes sense because these markets are influenced by factors that Bitcoin's dynamics do not adequately account for.

From a practical standpoint, this signal should be used as part of a broader risk management process. If the signal were used on its own as a warning, its importance would be overstated. The signal does not predict the next day's returns. Instead, it provides a small and relatively consistent shift toward caution as the structure of dependencies changes. This is a limited but reasonable use of the signal.

Future work should test the system over longer forecasting horizons, where the benefits of consistency may not be as evident. Future work should also consider a broader set of characteristics, including chain-linked indicators and macroeconomic variables, and compare the results with dependent indicators that capture extreme movements in joint variances more accurately than Pearson's moving correlation. A portfolio test that incorporates transaction costs will help make the results more realistic. These ideas are left for future work, and this thesis serves as a foundation for them.
